\section{Conclusion}
\label{sec:conclusion}

This paper evaluated whether the elementary mechanisms of a compound motor
fault can be recovered when their exact combination, their compound
context, and the operating condition are varied independently of the
training history. On a public multimodal induction-motor dataset whose
compound faults span mechanical and electrical modalities, a locked linear
constituent probe over mechanism-specific physical features showed that
strong closed-set discriminability (24-class accuracy 0.913, 95\% CI
$[0.885, 0.938]$) coexists with exact constituent-set recovery of only
0.083--0.157 on unseen compounds. The two results are not in tension; they
measure different abilities, and the gap between them is the paper's
subject.

Four conclusions follow from the controlled analyses. First, the failure
mode is incomplete decomposition rather than loss of information: at least
one true constituent was recovered in more than 92\% of compound runs, and
constituent recall (0.630--0.694) and all-constituent recovery under
recombination (up to 0.444) exceed input-independent prevalence references
(0.556 and 0.111, respectively), while several other partial-recovery
metrics do not---a distinction visible only because every metric was
reported against its constant-predictor floor. Second, compound-context
deprivation is measurably harmful chiefly when the operating condition is
simultaneously unseen: the recombination-versus-two-sided penalty under
condition shift was $+0.040$ micro-$F_1$ (95\% CI $[0.009, 0.073]$),
$+0.051$ constituent recall, and $+0.083$ all-constituent recovery, whereas
the corresponding seen-condition contrasts were not resolved. Third,
constituent evidence transfers across compound partners often but
directionally: within-pair detectability was near-perfect throughout, yet
cross-partner transfer ranged from perfect to a statistically supported
score-orientation reversal (Holm-adjusted $p \le 0.014$)---and the
demonstrated separability of distinct recordings of the same fault state
requires that reversal to be interpreted as context-dependent score
transfer rather than as proven physical signature inversion. Fourth,
absolute decomposition quality is an operating-point and model property:
precision-weighted thresholding raised exact-match recovery to
0.157--0.287 at reduced recall, and nonlinear classifier families reduced
false additions from ${\sim}0.68$ to ${\sim}0.26$ per run while the
qualitative generalization pattern persisted across all four families.

For practice, these results caution against reading closed-set accuracy as
deployment readiness, argue for calibrating constituent detectors with
explicit precision--recall cost structure, and indicate that commissioning
data should sample compound contexts and operating conditions jointly
rather than relying on either alone. The reference experiments of
Appendix~\ref{app:condition} reinforce the latter point: on this dataset
the standard leave-one-condition-out test is nearly saturated, while
single-condition training collapses accuracy, so robustness claims should
be evaluated in the extrapolating direction.

Several directions remain open. The dataset provides no replicated
same-state recordings that would separate fault, assembly, sensor-mounting,
and session effects, so the transfer analyses stop at carefully bounded
attribution; replicated acquisitions would convert those bounds into causal
statements. Representation or normalization strategies that promote
constituent-score consistency across compound partners are identified as
hypotheses, not established remedies. Finally, the constituent-level
protocol, prevalence references, and controls used here are directly
portable to the companion gearbox release and to other multimodal datasets,
where cross-dataset replication of the context-by-condition interaction
would test its generality.
